\documentclass[10pt,a4paper]{amsart}
\usepackage[margin=19mm]{geometry}
\usepackage{amsmath,amssymb,mathtools}
\usepackage[scr=rsfs,cal=euler]{mathalfa}
\usepackage{xfrac}
\usepackage[unicode,hidelinks,bookmarks=false]{hyperref}
\newcommand{\ie}{i.e.\ }
\newcommand{\cf}{cf.\ }
\newcommand{\resp}{respectively\ }
\newcommand{\Subsection}{\subsection}
\providecommand{\nobreakdash}{\nobreak\hbox{-}}
\setlength{\emergencystretch}{2em}
\multlinegap0pt
\usepackage{bm}
\usepackage{bbm}
\usepackage{dsfont}
\usepackage{xspace}
\usepackage{cancel}
\usepackage{mathtools}
\usepackage{amsthm}
\makeatletter
\@ifundefined{lemma}{\newtheorem{lemma}{Lemma}}{}
\@ifundefined{prop}{\newtheorem{prop}{Proposition}}{}
\makeatother
\title[Counting metacyclic fields]{Counting metacyclic fields}
\author{Maddie Allen, Justine Dell, Milad Fakhari, Kevin J. McGown, Chloe Stewart, Daniel Tedeschi}
\date{Source: arXiv:2608.09851v1 (CC BY 4.0)}
\begin{document}
\maketitle

\noindent\emph{Source and licence.} Counting metacyclic fields. Version arXiv:2608.09851v1 (CC BY 4.0), distributed under the Creative Commons Attribution 4.0 International licence, as recorded in the arXiv licence metadata for this exact version and shown on the abstract page https://arxiv.org/abs/2608.09851v1. Full rights notice: licenses/arxiv-2608.09851-CC-BY-4.0.txt.\par\smallskip

\noindent\textit{Editorial scope.} Counted unit: the Prop whose source label is \texttt{P:2} in the article (source0.tex lines 415--425, source label \texttt{P:2}), delivered together with the article's own complete original proof of it (source0.tex lines 428--467, one \texttt{proof} environment of 40 lines). No mathematics is written, repaired or re-derived: statement and proof are the article's own text line for line. Required context retained uncounted: L:1 (lines 183--199); L:2 (lines 214--231); L:sum (lines 244--261); L:constant (lines 306--309); L:sum2 (lines 331--341). Every definition, display and label of those lines is retained unchanged. Omitted from this package and not redistributed: 1--182, 200--213, 232--243, 262--305, 310--330, 342--414, 426--427, 468--530. Nothing inside a retained excerpt is omitted or altered: every retained body is delivered exactly as the article's lines are written, so no omission receipt of any kind accompanies this package. Citations: the retained excerpts carry no citation command at all, so no bibliography entry of the article is transcribed and no reference is redistributed. Reference adaptation: none was needed and none was made; every reference inside the delivered unit resolves inside it, so no unresolved reference or citation is produced. Printed slips kept literal: no printed word, symbol, hypothesis or number of the counted statement or of its proof was altered. Layout and class: the article's own class is \texttt{amsart} with packages including geometry, amsmath, amssymb, hyperref, xcolor, colonequals; the wrapper is the lane's pinned \texttt{amsart} editorial wrapper at 10pt on A4, which restates the theorem-like environments the retained text needs and adds the article's own macro definitions that the retained text needs (none), with extra wrapper packages (bm, bbm, dsfont, xspace, cancel, mathtools). The delivered numbering is the unit's own. Assets: no retained span contains a figure environment, a graphics-inclusion command, a PDF-page inclusion, a table or a TikZ picture; no image file is copied into this package, the package declares no asset dependency and no image file is redistributed with it. Licence: version arXiv:2608.09851v1 of this article is distributed under the Creative Commons Attribution 4.0 International licence, as recorded in the arXiv licence metadata for this exact version and shown on the abstract page https://arxiv.org/abs/2608.09851v1. A full rights notice is delivered at licenses/arxiv-2608.09851-CC-BY-4.0.txt. Characters: every retained excerpt is pure ASCII, so the delivered engine, XeLaTeX with its default Latin Modern text font, reports no missing character.
\begin{lemma}\label{L:1}
Notation as above.  The discriminants of $F$, $K$, $k$ are given by
$$
  \Delta_{F}=\Delta_k\, f^{\ell-1}
  \,,\quad
  \Delta_K=\Delta_k^{\ell}\, f^{(\ell-1)^2}
  \,,\quad
  \Delta_k=(-1)^{\frac{\ell-1}{2}}\ell^{\ell-2}
$$
where the associated conductor satisfies
\begin{align*}
f^{\ell-1}=\begin{cases}
\ell^2 R^{\ell-1} & D^{\ell-1}\not\equiv 1\pmod{\ell^2}\,,\\
R^{\ell-1} & D^{\ell-1}\equiv 1\pmod{\ell^2}\,.
\end{cases}
\end{align*}
\end{lemma}

\begin{lemma}\label{L:2}
Let $f$ be the conductor associated with a pure field of degree $\ell$.
Then
$$
m(f)=\begin{cases}
(\ell-1)^t\,, &e=\ell+1\,,\\ 
(\ell-1)^{u} a(v)\,, & e=2\,,\\
(\ell-1)^{u} a(v-1)\,, & e=0\,.
\end{cases}
$$
where
\begin{align*}
%u&=\#\{1\leq j\leq t\mid q_j^{\ell-1}\equiv 1\pmod{\ell^2}\}\\
%v&=\#\{1\leq j\leq t\mid q_j^{\ell-1}\not\equiv 1\pmod{\ell^2}\}\\
a(j)&=\frac{1}{\ell}\left((\ell-1)^j-(-1)^j\right)
\,,\quad j \geq -1\,.
\end{align*}
\end{lemma}

\begin{lemma}\label{L:sum}
We have
$$
\sum_{n \leq Y} \mu(n\ell)^2(\ell-1)^{\omega(n)}
\sim
B_\ell
Y(\log Y)^{\ell-2}
\,,
$$
$$
  B_\ell=
  e^{C_\ell}\,\frac{e^{-\gamma(\ell-1)}}{\Gamma(\ell-1)} \left( 1 + \frac{\ell-1}{\ell}\right)^{-1} \prod_{p \leq \ell-1} \left( 1 + \frac{\ell-1}{p} \right)
  \,,
$$
where $\gamma$ is the Euler--Mascheroni constant, $\omega(n)$ denotes the number of distinct prime divisors of $n$, and 
\[ C_{\ell}= (\ell-1)\left(M - \sum_{p\leq \ell-1} \frac{1}{p}\right) +\sum_{n=2}^\infty (-1)^{n+1}\frac{(\ell-1)^n P(\ell;n)}{n}\] for $M$ the Meissel--Mertens constant
\[ M:=\gamma+ \sum_p \left[\log\left(1 - \frac{1}{p} \right) + \frac{1}{p} \right].\]
\end{lemma}

\begin{lemma} \label{L:constant}
    We have 
    \[ B_{\ell} = \frac{1}{(\ell-2)!} \left( 1 + \frac{\ell-1}{\ell} \right)^{-1}\prod_{p} \left[ \left( 1 + \frac{\ell-1}{p} \right)\left( 1 - \frac{1}{p} \right)^{\ell-1} \right].\]
\end{lemma}

\begin{lemma}\label{L:sum2}
We have
$$
\sum_{n \leq Y}
%(-1)^{v(n)}
\mu(n\ell)^2(\ell-1)^{u(n)}
\ll_\ell
Y(\log Y)^{1-\frac{2}{\ell}}
\,.
$$
\end{lemma}

\begin{prop}\label{P:2}
Let $N_\ell^{(2)}(X)$ denote the number of $K\in\mathcal{M}_\ell$ 
with $|\Delta_K|\leq X$ such that $\ell\nmid D$ and $D^{\ell-1}\not\equiv 1\pmod{\ell^2}$.
We have
$$N_\ell^{(2)}(X)\sim A_\ell^{(2)} X^{\frac{1}{(\ell-1)^2}}(\log X)^{\ell-2}$$
where
$$
  A_\ell^{(2)}=B_{\ell}\,\ell^{-\frac{\ell^2-2}{(\ell-1)^2}-1}(\ell-1)^{-2(\ell-2).}
  \,.
$$
\end{prop}

\begin{proof}
Suppose $\ell\nmid D$
and
$D^{\ell-1}\not\equiv 1\pmod{\ell^2}$.
%Then $R = \ell q_1 \dots q_s$ for primes $q_i \neq \ell$.
By Lemma~\ref{L:1},
\begin{align*}
%f &= (1-\zeta_\ell)^{\ell+1}(q_1 \dots q_s)
%\,,\\
|\Delta_K| &= \ell^{\ell^2-2}(q_1 \dots q_s)^{(\ell-1)^2}
\,,
\end{align*}
and hence $|\Delta_K|\leq X$ implies
$$
  q_1\dots q_s\leq \left(\frac{X}{\ell^{\ell^2-2}}\right)^{\frac{1}{(\ell-1)^2}}=:Y
  \,.
$$
As above, counting that number of metacyclic fields with $|\Delta_K|\leq X$ in this case
amounts to counting the number of squarefree numbers $n$ not divisible by $\ell$
such that $n\leq Y$ with the multiplicity
--- given by Lemma~\ref{L:2} --- equal to $(\ell-1)^{u(n)}a(v(n))$. 
Therefore, the contribution in this case is
\begin{align*}
  N_\ell^{(2)}(X)&=\sum_{n\leq Y} \mu(n\ell)^2(\ell-1)^{u(n)}a(v(n))
\\
% &=
% \sum_{n\leq Y} \mu(n\ell)^2(\ell-1)^{u(n)}
% \frac{1}{\ell}\left((\ell-1)^{v(n)}-(-1)^{v(n)}\right)
% \\
&=
\frac{1}{\ell}
\sum_{n\leq Y} \mu(n\ell)^2(\ell-1)^{\omega(n)}
-
\frac{1}{\ell}
\sum_{n\leq Y} \mu(n\ell)^2(\ell-1)^{u(n)}
(-1)^{v(n)}
\,.
\end{align*}
Applying Lemmas~\ref{L:sum}, \ref{L:constant}, and~\ref{L:sum2} proves the result.
\end{proof}

\end{document}
